\documentclass[11pt]{article}
\usepackage[margin=1in]{geometry}
\usepackage{amsmath,amssymb,amsthm}
\usepackage{hyperref}
\usepackage{enumitem}

\newtheorem{theorem}{Theorem}
\newtheorem{lemma}[theorem]{Lemma}
\newtheorem{corollary}[theorem]{Corollary}
\theoremstyle{definition}
\newtheorem{definition}[theorem]{Definition}
\theoremstyle{remark}
\newtheorem{remark}[theorem]{Remark}

\usepackage{xurl}
\let\oldus\_
\renewcommand{\_}{\oldus\allowbreak}
\newcommand{\DM}{\operatorname{DM}}

\title{A disproof of Conjecture 00000002591\\[2pt]
\large No finite poset with at least three elements has $2^{|P|}$ Dedekind--MacNeille cuts}
\date{}

\begin{document}
\sloppy
\maketitle

\section{The conjecture}

Conjecture 00000002591 reads:

\begin{quote}
\emph{Definition:} The Dedekind--MacNeille completion $\DM(P)$ is the smallest complete extension of a
poset $P$ preserving all existing joins and meets. \emph{Conjecture:} The size law of MacNeille
completions: $|\DM(P)| \le 2^{|P|}$ with the upper bound attained by unbounded dense antichains; the
isomorphism types of $P$ attaining the bound are superpolynomial in number, with the minimal generator
characterization by three-chain cross products.
\end{quote}

It is a conjunction of four clauses:
\begin{enumerate}[label=(\roman*)]
\item $|\DM(P)| \le 2^{|P|}$;
\item the bound is attained by (unbounded dense) antichains;
\item the number $E(n)$ of isomorphism types of $n$-element posets $P$ with $|\DM(P)| = 2^n$ is
  superpolynomial in $n$;
\item a ``minimal generator characterization by three-chain cross products''.
\end{enumerate}
We prove (i) for finite posets (the Lean theorem \texttt{card\_cuts\_le} assumes \texttt{Fintype}; the
same injection of cuts into subsets gives the inequality for arbitrary cardinals, but that general
statement is not formalized), and we refute (iii) and (ii) for finite posets: $E(n) = 0$ for every
$n \ge 3$.
Clause (iv) is too vague to formalize and is not addressed; refuting (iii) already refutes the
conjunction.

\paragraph{Readings of (iii) that are refuted.} $E(n)$ is counted either up to isomorphism or as labelled
partial orders on $\{0,\dots,n-1\}$; ``attaining the bound'' is read either as $|\DM(P)| = 2^n$ or as
$|\DM(P)| \ge 2^n$ (by (i) these are the same); ``superpolynomial'' is read either strongly (for every
$k$, $E(n) > n^k$ for all large $n$) or weakly (for every $k$ and $C$, $E(n) > C n^k$ for infinitely many
$n$, i.e.\ $E$ is not $O(n^k)$ for any $k$). All of these are refuted.

\paragraph{Not refuted / not formalized.} Clause (ii) for infinite posets is not formalized. (In prose:
an infinite antichain of cardinality $\kappa$ has $\kappa + 2 = \kappa < 2^\kappa$ cuts, so infinite
antichains do not attain the bound either; infinite dense linear orders such as $\mathbb{Q}$ do attain
$2^{\aleph_0}$, but they are not antichains.) Clause (iv) is not addressed.

\section{Definitions and conventions}

For $A \subseteq P$ let $A^u$ be the set of upper bounds of $A$ and $A^\ell$ the set of lower bounds.
Following the retrieved Wikipedia article on the Dedekind--MacNeille completion
(\url{https://en.wikipedia.org/wiki/Dedekind%E2%80%93MacNeille_completion}), $\DM(P)$ ``consists of all
subsets $A$ for which $(A^u)^\ell = A$, ordered by inclusion''; equivalently a \emph{cut} is ``a pair of
sets $(A, B)$ for which $A^u = B$ and $A = B^\ell$''.

In Lean we use Mathlib's \texttt{DedekindCut P}, which is \texttt{Concept P P (. <= .)}: pairs
(\texttt{left}, \texttt{right}) with \texttt{upperBounds left = right} and \texttt{lowerBounds right =
left}. This is exactly the pair form of a cut; a cut is determined by its left set
(\texttt{Concept.extent\_injective}), and a set $s$ is a left set iff
$\texttt{lowerBounds}(\texttt{upperBounds}\ s) = s$ (Lean: \texttt{exists\_cut\_iff}). So
$|\DM(P)|$ is \texttt{Nat.card (DedekindCut P)}.

\emph{Empty-set convention.} No extra elements are adjoined. The empty set is a left set iff $P$ has no
least element; $P$ itself is always a left set. For example, an antichain with $n \ge 2$ elements has the
$n+2$ cuts $\emptyset$, the $n$ singletons and $P$, matching the Wikipedia description of the completion
of an antichain as ``$S$ itself together with two additional elements''. The argument below uses only
subsets with one or two elements, so the conclusion does not change if $\emptyset$ and $P$ are always
counted as cuts.

\emph{Counting.} Isomorphism types of $n$-element posets are represented by partial orders on
\texttt{Fin n}. \texttt{Attaining n} is the type of partial orders $Q$ on \texttt{Fin n} with $2^n \le
|\texttt{DedekindCut}|$ for $Q$, and $E(n)$ is the cardinality of its quotient by order isomorphism:
\texttt{E n = Nat.card (Quot Iso)}. Both types are finite (Lean instances), so $E(n)$ is a genuine count.

\section{Proof}

\begin{lemma}[\texttt{no\_three\_of\_cuts}]
If every singleton and every two-element subset of $P$ is the left set of a cut, then $P$ has no three
pairwise distinct elements.
\end{lemma}

\begin{proof}
For $u \in P$, $\{u\}^u = {\uparrow}u$ and $({\uparrow}u)^\ell = {\downarrow}u$, so $\{u\}$ being closed
gives ${\downarrow}u = \{u\}$ for all $u$. Let $x, y, z$ be pairwise distinct. If $u$ is an upper bound of
$\{x, y\}$ then $x, y \in {\downarrow}u = \{u\}$, so $x = u = y$, a contradiction; hence $\{x,y\}^u =
\emptyset$ and $(\{x,y\}^u)^\ell = \emptyset^\ell = P$. Since $\{x,y\}$ is closed, $P = \{x, y\}$, so
$z \in \{x, y\}$, a contradiction.
\end{proof}

\begin{lemma}[\texttt{exists\_proper\_not\_cut}]
If $|P| \ge 3$, some nonempty proper subset of $P$ (of size one or two) is not the left set of a cut.
\end{lemma}

\begin{proof}
Pick three distinct elements. If some singleton $\{u\}$ is not a left set, it is nonempty and proper
(it misses one of the other elements). Otherwise, by the previous lemma some pair $\{u, v\}$ is not a
left set; it is proper because it cannot contain three distinct elements.
\end{proof}

\begin{theorem}[\texttt{card\_cuts\_le}, \texttt{card\_cuts\_lt}]
For a finite poset $P$, $|\DM(P)| \le 2^{|P|}$; if $|P| \ge 3$ then $|\DM(P)| < 2^{|P|}$.
\end{theorem}

\begin{proof}
The map from cuts to their left sets is injective into the power set, which has $2^{|P|}$ elements. If
$|P| \ge 3$, the previous lemma gives a subset outside its range, so the inequality is strict.
\end{proof}

In particular no antichain with at least three elements attains the bound, which refutes clause (ii) for
finite antichains.

\begin{theorem}[\texttt{not\_superpolynomial}]
$E(n) = 0$ for all $n \ge 3$. Consequently it is false that for every $k$, $E(n) > n^k$ eventually, and
it is false that for every $k, C$, $E(n) > C n^k$ infinitely often.
\end{theorem}

\begin{proof}
By the previous theorem, \texttt{Attaining n} is empty for $n \ge 3$ (\texttt{attaining\_isEmpty}), so
both the labelled count (\texttt{card\_attaining\_eq\_zero}) and $E(n)$ are $0$. Taking $k = 0$ (and
$C = 0$), either growth statement would give some $n \ge 3$ with $E(n) > 0$.
\end{proof}

\begin{remark}[non-vacuity, \texttt{antichain2\_attains}, \texttt{E\_two\_pos}]
The bound is sharp at $n = 2$: the two-element antichain on \texttt{Fin 2} has all four subsets as left
sets of cuts, so $E(2) > 0$. This checks that the definitions of \texttt{Attaining} and $E$ are not
vacuous. A brute-force check also confirms that no poset on $3$ or $4$ points has $2^n$ cuts and that
antichains on $3$ and $4$ points have $n + 2$ cuts.
\end{remark}

\section{Lean formalization}

File \texttt{lean/Conjecture2591/Basic.lean} (namespace \texttt{Conjecture2591}, 206 lines, only import
\texttt{Mathlib}). Main theorem (ASCII rendering):
\begin{verbatim}
theorem not_superpolynomial :
    (forall n, 3 <= n -> E n = 0) /\
    not (forall k : Nat, forall^f n in atTop, n ^ k < E n) /\
    not (forall k C : Nat, exists^f n in atTop, C * n ^ k < E n)
\end{verbatim}
(\texttt{forall\^{}f} and \texttt{exists\^{}f} stand for Lean's \texttt{Filter.Eventually} and
\texttt{Filter.Frequently}.) Supporting results: \texttt{exists\_cut\_iff}, \texttt{no\_three\_of\_cuts},
\texttt{exists\_proper\_not\_cut}, \texttt{card\_cuts\_le}, \texttt{card\_cuts\_lt},
\texttt{attaining\_isEmpty}, \texttt{card\_attaining\_eq\_zero}, \texttt{E\_eq\_zero},
\texttt{antichain2\_attains}, \texttt{E\_two\_pos}.

\paragraph{Axiom audit.} \texttt{\#print axioms} for \texttt{not\_superpolynomial},
\texttt{card\_cuts\_lt} and \texttt{E\_two\_pos} reports only \texttt{propext}, \texttt{Classical.choice},
\texttt{Quot.sound}. There is no \texttt{sorry} and no \texttt{native\_decide}.

\end{document}
